\documentclass[10pt,a4paper]{amsart}
\usepackage{amsmath,amssymb,amsthm,graphicx}
% Retained minimal preamble for the Stacks Project chapter varieties.tex.
% Verbatim excerpt of the chapter preamble: Xy-pic setup (native preamble lines 14-18), the native theorem
% declarations with the mathematical macros (native preamble lines 154-217), and the project own no-op
% definitions of the reference/slogan/history display environments (native wrapper.tex lines 1-3) that the
% counted statement uses for its slogan. Copied mechanically; nothing added, removed or redefined.
% For commutative diagrams we use Xy-pic
\usepackage[all]{xy}

% We use 2cell for 2-commutative diagrams.
\xyoption{2cell}
% Theorem environments.
%
\theoremstyle{plain}
\newtheorem{theorem}[subsection]{Theorem}
\newtheorem{proposition}[subsection]{Proposition}
\newtheorem{lemma}[subsection]{Lemma}

\theoremstyle{definition}
\newtheorem{definition}[subsection]{Definition}
\newtheorem{example}[subsection]{Example}
\newtheorem{exercise}[subsection]{Exercise}
\newtheorem{situation}[subsection]{Situation}

\theoremstyle{remark}
\newtheorem{remark}[subsection]{Remark}
\newtheorem{remarks}[subsection]{Remarks}

\numberwithin{equation}{subsection}

% Macros
%
\def\lim{\mathop{\mathrm{lim}}\nolimits}
\def\colim{\mathop{\mathrm{colim}}\nolimits}
\def\Spec{\mathop{\mathrm{Spec}}}
\def\Hom{\mathop{\mathrm{Hom}}\nolimits}
\def\Ext{\mathop{\mathrm{Ext}}\nolimits}
\def\SheafHom{\mathop{\mathcal{H}\!\mathit{om}}\nolimits}
\def\SheafExt{\mathop{\mathcal{E}\!\mathit{xt}}\nolimits}
\def\Sch{\mathit{Sch}}
\def\Mor{\mathop{\mathrm{Mor}}\nolimits}
\def\Ob{\mathop{\mathrm{Ob}}\nolimits}
\def\Sh{\mathop{\mathit{Sh}}\nolimits}
\def\NL{\mathop{N\!L}\nolimits}
\def\CH{\mathop{\mathrm{CH}}\nolimits}
\def\proetale{{pro\text{-}\acute{e}tale}}
\def\etale{{\acute{e}tale}}
\def\QCoh{\mathit{QCoh}}
\def\Ker{\mathop{\mathrm{Ker}}}
\def\Im{\mathop{\mathrm{Im}}}
\def\Coker{\mathop{\mathrm{Coker}}}
\def\Coim{\mathop{\mathrm{Coim}}}

% Boxtimes
%
\DeclareMathSymbol{\boxtimes}{\mathbin}{AMSa}{"02}

%
% Macros for moduli stacks/spaces
%
\def\QCohstack{\mathcal{QC}\!\mathit{oh}}
\def\Cohstack{\mathcal{C}\!\mathit{oh}}
\def\Spacesstack{\mathcal{S}\!\mathit{paces}}
\def\Quotfunctor{\mathrm{Quot}}
\def\Hilbfunctor{\mathrm{Hilb}}
\def\Curvesstack{\mathcal{C}\!\mathit{urves}}
\def\Polarizedstack{\mathcal{P}\!\mathit{olarized}}
\def\Complexesstack{\mathcal{C}\!\mathit{omplexes}}
% \Pic is the operator that assigns to X its picard group, usage \Pic(X)
% \Picardstack_{X/B} denotes the Picard stack of X over B
% \Picardfunctor_{X/B} denotes the Picard functor of X over B
\def\Pic{\mathop{\mathrm{Pic}}\nolimits}
\def\Picardstack{\mathcal{P}\!\mathit{ic}}
\def\Picardfunctor{\mathrm{Pic}}
\def\Deformationcategory{\mathcal{D}\!\mathit{ef}}
\newenvironment{reference}{}{}
\newenvironment{slogan}{}{}
\newenvironment{history}{}{}
\usepackage{cleveref}
\crefname{originallemma}{original lemma}{original lemmas}
\Crefname{originallemma}{Original lemma}{Original lemmas}
\makeatletter
\newlabel{lemma-dim-1-projective-completion}{{0BXV}{source}{}{sourcepointer-1}{}}
\newlabel{lemma-dim-1-projective-completion@cref}{{[originallemma][0][]0BXV}{[1][0][]source}{}{}{}}
\makeatother
\crefname{originallemma}{original lemma}{original lemmas}
\Crefname{originallemma}{Original lemma}{Original lemmas}
\makeatletter
\newlabel{lemma-relative-normalization-finite}{{0BXS}{source}{}{sourcepointer-2}{}}
\newlabel{lemma-relative-normalization-finite@cref}{{[originallemma][0][]0BXS}{[1][0][]source}{}{}{}}
\makeatother
\crefname{originalreference}{original reference}{original references}
\Crefname{originalreference}{Original reference}{Original references}
\makeatletter
\newlabel{morphisms-lemma-finite-projective}{{0B3I}{source}{}{sourcepointer-3}{}}
\newlabel{morphisms-lemma-finite-projective@cref}{{[originalreference][0][]0B3I}{[1][0][]source}{}{}{}}
\makeatother
\crefname{originalreference}{original reference}{original references}
\Crefname{originalreference}{Original reference}{Original references}
\makeatletter
\newlabel{morphisms-lemma-projective-over-quasi-projective-is-H-projective}{{087S}{source}{}{sourcepointer-4}{}}
\newlabel{morphisms-lemma-projective-over-quasi-projective-is-H-projective@cref}{{[originalreference][0][]087S}{[1][0][]source}{}{}{}}
\makeatother
\crefname{originalreference}{original reference}{original references}
\Crefname{originalreference}{Original reference}{Original references}
\makeatletter
\newlabel{morphisms-lemma-H-projective-composition}{{01WE}{source}{}{sourcepointer-5}{}}
\newlabel{morphisms-lemma-H-projective-composition@cref}{{[originalreference][0][]01WE}{[1][0][]source}{}{}{}}
\makeatother
\crefname{originalreference}{original reference}{original references}
\Crefname{originalreference}{Original reference}{Original references}
\makeatletter
\newlabel{morphisms-lemma-relative-normalization-normal-codim-1}{{0BXB}{source}{}{sourcepointer-6}{}}
\newlabel{morphisms-lemma-relative-normalization-normal-codim-1@cref}{{[originalreference][0][]0BXB}{[1][0][]source}{}{}{}}
\makeatother
\title{A reduced separated scheme of finite type of dimension at most one admits an H-projective completion with discrete valuation ring boundary}
\author{The Stacks Project Authors}
\date{stacks-project:varieties.tex:local-snapshot-20261011; GNU FDL 1.2 or later}
\begin{document}
\maketitle
\noindent\textit{Editorial note.} All delivered mathematical text is the author's own, copied verbatim from the pinned Stacks Project chapter file varieties.tex; no statement or proof step was rewritten, reordered or removed.\par\medskip
\noindent\textit{Editorial note.} The counted span cites six cross-chapter labels, each recorded in assembly.external\_references with its canonical tag as a navigation pointer.\par\medskip
\noindent\textit{Editorial note.} The proof cites no bibliography item inside the counted span.\par\medskip
\noindent\textit{Allowed standard prerequisites (editorial).}\begin{itemize}
\item Projective completion of a separated scheme of finite type of dimension at most one. Stacks Project, Varieties, Lemma lemma-dim-1-projective-completion (tag 0A26). Produces the initial immersion X -> Xbar that is then normalized.
\item Relative normalization is finite and normal schemes are discrete valuation rings in codimension one. Stacks Project, Morphisms, Lemma morphisms-lemma-relative-normalization-finite (tag 0BXL) and Lemma morphisms-lemma-relative-normalization-normal-codim-1 (tag 0BXN). The first gives finiteness of the normalization and the second identifies the boundary local rings as discrete valuation rings.
\end{itemize}\medskip
\ifdefined\hypertarget\hypertarget{sourcepointer-1}{}\fi
\noindent\textit{Source reference.} Original lemma 0BXV: sources/stacks-chow/tags.txt, source lines 14993--14993. Cross-chapter standard prerequisite cited by the author inside the counted span; it is used as a pointer and its own proof is not part of this item.\par\medskip
\ifdefined\hypertarget\hypertarget{sourcepointer-2}{}\fi
\noindent\textit{Source reference.} Original lemma 0BXS: sources/stacks-chow/tags.txt, source lines 14990--14990. Cross-chapter standard prerequisite cited by the author inside the counted span; it is used as a pointer and its own proof is not part of this item.\par\medskip
\ifdefined\hypertarget\hypertarget{sourcepointer-3}{}\fi
\noindent\textit{Source reference.} Original reference 0B3I: sources/stacks-chow/tags.txt, source lines 13966--13966. Cross-chapter standard prerequisite cited by the author inside the counted span; it is used as a pointer and its own proof is not part of this item.\par\medskip
\ifdefined\hypertarget\hypertarget{sourcepointer-4}{}\fi
\noindent\textit{Source reference.} Original reference 087S: sources/stacks-chow/tags.txt, source lines 10437--10437. Cross-chapter standard prerequisite cited by the author inside the counted span; it is used as a pointer and its own proof is not part of this item.\par\medskip
\ifdefined\hypertarget\hypertarget{sourcepointer-5}{}\fi
\noindent\textit{Source reference.} Original reference 01WE: sources/stacks-chow/tags.txt, source lines 2479--2479. Cross-chapter standard prerequisite cited by the author inside the counted span; it is used as a pointer and its own proof is not part of this item.\par\medskip
\ifdefined\hypertarget\hypertarget{sourcepointer-6}{}\fi
\noindent\textit{Source reference.} Original reference 0BXB: sources/stacks-chow/tags.txt, source lines 14974--14974. Cross-chapter standard prerequisite cited by the author inside the counted span; it is used as a pointer and its own proof is not part of this item.\par\medskip
\begin{lemma}
\label{lemma-reduced-dim-1-projective-completion}
Let $X$ be a separated scheme of finite type over $k$.
If $X$ is reduced and $\dim(X) \leq 1$, then there exists
an open immersion $j : X \to \overline{X}$ such that
\begin{enumerate}
\item $\overline{X}$ is H-projective over $k$, i.e., $\overline{X}$
is a closed subscheme of $\mathbf{P}^d_k$ for some $d$,
\item $j(X) \subset \overline{X}$ is dense and scheme
theoretically dense,
\item $\overline{X} \setminus X = \{x_1, \ldots, x_n\}$
for some closed points $x_i \in \overline{X}$,
\item the local rings $\mathcal{O}_{\overline{X}, x_i}$
are discrete valuation rings for $i = 1, \ldots, n$.
\end{enumerate}
\end{lemma}

\begin{proof}
Let $j : X \to \overline{X}$ be as in
Lemma \ref{lemma-dim-1-projective-completion}.
Consider the normalization $X'$ of $\overline{X}$ in
$X$. By Lemma \ref{lemma-relative-normalization-finite}
the morphism $X' \to \overline{X}$ is finite.
By Morphisms, Lemma \ref{morphisms-lemma-finite-projective}
$X' \to \overline{X}$ is projective. By Morphisms, Lemma
\ref{morphisms-lemma-projective-over-quasi-projective-is-H-projective}
we see that $X' \to \overline{X}$ is H-projective.
By Morphisms, Lemma \ref{morphisms-lemma-H-projective-composition}
we see that $X' \to \Spec(k)$ is H-projective.
Let $\{x'_1, \ldots, x'_m\} \subset X'$ be the inverse
image of $\{x_1, \ldots, x_n\} = \overline{X} \setminus X$.
Then $\dim(\mathcal{O}_{X', x'_i}) = 1$ for all $1 \leq i \leq m$.
Hence the local rings $\mathcal{O}_{X', x'}$
are discrete valuation rings by
Morphisms, Lemma
\ref{morphisms-lemma-relative-normalization-normal-codim-1}.
Then $X \to X'$ and $\{x'_1, \ldots, x'_m\}$ is as desired.
\end{proof}

\end{document}
